\chapter{Comparative Statics Analysis}

A central question in economics is how an endogenous variable moves
with an exogenous parameter. Take the firm's indirect single-product
profit maximization:
\[\max _{y\ge0}\;py-c(\w,y).\]
Holding the input price vector \(\w\) fixed, how does the optimal
supply \(y^*\) respond to the output price \(p\)?

More generally, given an objective \(F: X\times \Theta \to \RR\) with
\(X,\Theta\subset \RR\), we ask how the maximizer
\[x^*(\theta) = \arg\max_{x\in X}F(x,\theta)\]
moves with the parameter \(\theta\). The two competing approaches we
will develop are the \emph{classical} (calculus-based, using the
implicit function theorem) and the \emph{monotone} or \emph{lattice}
approach (order-theoretic, requiring only complementarity). The latter
is more robust — it works without differentiability, without
concavity, and even when \(x^*(\theta)\) is set-valued.

\section{Univariate Comparative Statics}\label{univariate-comparative-statics}

The classical approach uses the implicit function theorem. It requires
two sets of assumptions:

\begin{itemize}
\item
  \(F(\cdot,\cdot )\) is twice continuously differentiable.
\item
  The maximizer \(x^*(\theta)\) is unique and is characterized by the
  first-order condition.
\end{itemize}

\rmkb{
A set of sufficient conditions for the second assumption is needed:

\begin{itemize}
\item
  The choice set \(X\) is convex.
\item
  \(F(\cdot,\theta)\) is strictly concave in \(x\).
\item
  The solution is interior.
\end{itemize}
}

Under the two sets of assumptions, the first-order condition is:
\[F_x(x(\theta),\theta) = 0.\]

Differentiating on both sides with respect to \(\theta\), we have: 
\[\ba
& F_{xx}(x(\theta),\theta)\cdot x'(\theta)+F_{x\theta}(x(\theta),\theta)=0\\
\too & x'(\theta) = -\frac{F_{x\theta}(x(\theta),\theta)}{F_{xx}(x(\theta),\theta)}
\ea\]

If \(F\pr{\cdot,\theta}\) is strictly concave in \(x\), then
\(F_{xx}\pr{\cdot,\cdot}<0\) (assuming no reflection point), so
\(x\pr{\cdot}\) is strictly increasing in \(\theta\) if
\(F_{x\theta}\pr{\cdot,\cdot}>0\).

The classical method's strength is that it delivers a closed-form
expression for \(x'(\theta)\), which is useful for quantitative work.
The flip side is that the same machinery brings real drawbacks:
\begin{itemize}
  \item
    \textit{Technical:} the prerequisites (twice differentiability,
    strict concavity, interior solution) are strong, and the algebra
    can be tedious.
  \item
    \textit{Substantive:} strict concavity in \(x\) is a particularly
    awkward assumption, because concavity is not invariant under monotone
    transformations — but the comparative-statics \emph{direction} of
    \(x^*(\theta)\) is. Asking for concavity to detect a direction that
    does not depend on concavity is overkill.
  \end{itemize}

\ex{
Taking \(\w\) as given, the profit maximization problem is given by
\[\max_{y\ge 0} py-c(y)\] 

We use the classical approach to determine
whether (and when) the firm's supply curve is (weakly)
upward sloping.

Notice that the choice set \(Y = [0,+\infty)\) is convex, and when
\(c(\cdot)\) is strictly convex, the objective function is strictly
concave. Assuming interiority, the first-order condition is given by:
\[p = c'\pr{y\pr{p}}\]
Differentiating on both sides with respect to \(p\), we get
\[y'\pr{p} =\dfrac{1}{c''\pr{y\pr{p}}}\]

Thus, \(c\pr{\cdot}\) being
strictly convex is a sufficient condition for the supply curve being upward sloping. However, the requirement of
\(c\pr{\cdot}\) being strictly convex is not always a reasonable
assumption. Moreover, this assumption is not necessary. Recall that we
have \emph{Law of Supply}, which states that the firm's
supply curve is weakly upward sloping without any other assumption.
}


The diagnosis: differentiability and concavity of \(F(\cdot,\theta)\)
are not what comparative statics fundamentally needs. The real driver
is \(F_{x\theta}\ge 0\) — the \emph{complementarity} between \(x\) and
\(\theta\). The natural question is whether this cross-partial
condition (suitably generalized) is by itself enough for
\(x^*(\theta)\) to be increasing. The remainder of this section
develops a discrete, order-theoretic analogue that does not require
\(F\) to be differentiable or concave.

\defn{Increasing Differences}{
Let
\(F:X\times \Theta\to \RR\), where \(X,\Theta\subset \RR\). We say that
\(F(\cdot,\cdot)\) has \emph{increasing differences} in \((x,\theta)\)
if for any \(x,x'\in X\) and \(\theta,\theta'\in \Theta\) such that
\(x'>x\) and \(\theta'>\theta\), we have
\[F\pr{x',\theta'}-F\pr{x,\theta'}\ge F\pr{x',\theta}-F\pr{x,\theta}\]
If the inequalities are strict for all such \(x,x'\in X\) and
\(\theta,\theta'\in \Theta\), \(F\pr{\cdot,\cdot}\) has \emph{strictly
increasing differences} in \(\pr{x,\theta}\).
}

Intuitively, increasing differences says that the \emph{marginal
value} of \(x\) is larger at higher \(\theta\) — \(x\) and \(\theta\)
are \emph{complements}. Crucially, the condition is stated in terms of
discrete differences, so it makes sense even when \(F\) is not
differentiable, and it is invariant under monotone transformations of
\(F\) (unlike concavity).

\prop{Increasing Differences for Smooth Functions}{
Suppose
\(X = \br{\underline x,\bar x}\) and
\(\Theta = \br{\underline \theta,\bar \theta}\), where
\(X,\Theta\subset \RR\).

\begin{enumerate}
\item
  If \(F\pr{\cdot,\cdot}\) is continuously differentiable in both \(x\)
  and \(\theta\), \(F\pr{\cdot,\cdot}\) has increasing
  differences in \(\pr{x,\theta}\) if and only if either of the two conditions holds:

  \begin{itemize}
  \item
    \(F_x(x,\cdot)\) is non-decreasing in \(\theta\) for all \(x\).
  \item
    \(F_{\theta}\pr{\cdot,\theta}\) is non-decreasing in \(x\) for all
    \(\theta\).
  \end{itemize}
\item
  If \(F\pr{\cdot,\cdot}\) is twice continuously differentiable in both
  \(x\) and \(\theta\), \(F\pr{\cdot,\cdot}\) has increasing
  differences in \(\pr{x,\theta}\) if and only if
  \(F_{x\theta}\pr{\cdot,\cdot}\ge 0\) for all \(\pr{x,\theta}\).
\end{enumerate}
}

One last piece of plumbing. When \(F(\cdot,\theta)\) is not strictly
concave, \(x^*(\theta)\) is a \emph{set}, not a single number. We need
a way to say that one set lies ``below'' another. Two natural orders
on sets serve this purpose:

\defn{Comparison of Sets}{
For any two sets \(A\) and
\(B\), we say that:

\begin{itemize}
\item
  \(A\le B\) \emph{in the strong set order} if for any \(a\in A\) and
  \(b\in B\), we have \(\min \brace{a,b}\in A\) and
  \(\max\brace{a,b}\in B\).
\item
  \(A\le B\) \emph{pointwise} if for any \(a\in A\) and \(b\in B\), we
  have \(a\le b\).
\end{itemize}
}

Intuitively, the \emph{strong set order} allows \(A\) and \(B\) to
overlap (the overlap forms a shared region), but outside that overlap
every \(A\)-element lies below every \(B\)-element. The
\emph{pointwise} order is stricter: every element of \(B\) must lie at
or above \emph{every} element of \(A\), regardless of overlap.

\thmp{Univariate Topkis' Theorem}{
Let
\(F:X\times \Theta\to \RR\), where \(X,\Theta\subset \RR\), and
\(x^*\pr{\theta} = \arg\max_{x\in X}F\pr{x,\theta}\). Then for any
\(\theta'>\theta\),

\begin{enumerate}
\item
  If \(F\pr{\cdot,\cdot}\) has increasing differences in
  \(\pr{x,\theta}\), then \(x^*\pr{\theta}\le x^*\pr{\theta'}\) in the
  strong set order.
\item
  If \(F\pr{\cdot,\cdot}\) has strictly increasing differences in
  \(\pr{x,\theta}\), then \(x^*\pr{\theta}\le x^*\pr{\theta'}\)
  pointwise.
\end{enumerate}
}{
Take any \(x\in x^*\pr{\theta}\) and
\(x'\in x^*\pr{\theta'}\). Suppose \(x>x'\), then by revealed
preference,
\[\ba
F\pr{x,\theta}\ge F\pr{x',\theta}\\
F\pr{x',\theta'}\ge F\pr{x,\theta'}
\ea\]
By increasing differences,
\[F\pr{x,\theta}-F\pr{x',\theta}\le F\pr{x,\theta'}-F\pr{x',\theta'}\]
Jointly we have:
\[\ba
& 0\le F\pr{x,\theta}-F\pr{x',\theta}\le F\pr{x,\theta'}-F\pr{x',\theta'}\le 0\\
\too & \bc
F\pr{x,\theta} = F\pr{x',\theta}\\
F\pr{x,\theta'} = F\pr{x',\theta'}
\ec\too 
\bc
x\in x^*\pr{\theta'}\\
x'\in x^*\pr{\theta}
\ec
\ea\]

If \(F\pr{\cdot,\cdot}\) has strictly increasing differences in
\(\pr{x,\theta}\), then the combined inequality cannot hold, so it must
be that \(x^*\pr{\theta}\le x^*\pr{\theta'}\) pointwise. In particular, if \(F\pr{\cdot,\cdot}\) is twice continuously
differentiable in both \(x\) and \(\theta\) and \(x^*\pr{\theta}\) is
single-valued, then \(F_{x\theta}\pr{\cdot,\cdot}>0\) is sufficient for
\(x^*\pr{\theta}\) to be (weakly) increasing in \(\theta\).
}

\ex{
Again consider the
firm's profit maximization problem. How would the
firm's supply (correspondence) change with the output
price \(p\)?

The objective function is \(F\pr{y,p} = py-c\pr{\w,y}\).
\(Y = [0,+\infty)\) and \(P = [0,+\infty)\) are both intervals in
\(\RR\). Moreover, \(F_p\pr{y,p} = y\), which is strictly increasing in
\(y\), so \(F\pr{\cdot,\cdot}\) has strictly increasing differences. Therefore, the firm's supply increases with the output price $p$ pointwise.
}

\ex{
Consider a monopolist that faces a
downward sloping demand curve \(Q^D\pr{p}\). Now suppose the government
levies a unit tax \(t\) on the firm. How would the before-tax price
\(p\) received by the firm change with the unit tax \(t\)?

The objective function is
\(F\pr{p,t} = \pr{p-t}Q^D\pr{p}-c\pr{Q^D\pr{p}}\). \(P = [0,+\infty)\)
and \(T = [0,+\infty)\) are both intervals in \(\RR\). Moreover,
\(F_t\pr{p,t} = -Q^D\pr{p}\), which is strictly increasing in \(p\)
(since the demand curve is downward sloping), so \(F\pr{\cdot,\cdot}\)
has strictly increasing differences in \(\pr{p,t}\). Consequently, the
firm's before-tax price \(p\) increases with \(t\)
pointwise.
}

Notice that whether \(x^*\pr{\theta}\)
increases/decreases with the parameter \(\theta\) is an \emph{ordinal}
property, while (strictly) increasing differences is still a
\emph{cardinal} property. Indeed, we know from the discussion on
consumer theory that \(\max_xF\pr{x,\theta}\) and
\(\max_{x}G\pr{x,\theta}\) have the same set of maximizers if
\(G = \varphi\circ F\) for \(\varphi\pr{\cdot}\) strictly increasing.
Nevertheless, \(G\) having increasing differences in
\(\pr{x,\theta}\) does not necessarily mean \(F\) having (strictly)
increasing differences in \(\pr{x,\theta}\). In other words, the
requirement of (strictly) increasing differences is still too strong for
monotone comparative statics. For our purpose, if we can find a positive and monotonic transformation $\varphi$ such that \(G =\varphi \circ F\) has increasing differences or strictly increasing differences in \(\pr{x,\theta}\), then we know
\(x^*\pr{\theta}\) increases with \(\theta\) in the strong set order or pointwise.


\ex{
Consider the effects of an increase in
the market size on monopoly quantity (and monopoly price). Each consumer
in the market has an identical inverse function given by \(p^D\pr{q}\).
Suppose the number of consumers \(N\) is exogenously given, and that the
firm's cost function is \(c\pr{Q}\), where \(Q = Nq\) is
the total quantity sold (i.e., the number of consumers times per unit
purchase). Discuss how the firm's cost function
\(c\pr{\cdot}\) would affect the optimal \emph{per-consumer} quantity
\(q^*\pr{N}\) as the number of consumers \(N\) increases.
}

\sol{
The objective function is
\[F\pr{q,N} = N\cdot p^D\pr{q}\cdot q - c\pr{Nq}.\]

Notice that it is hard to check increasing differences of \(F\pr{\cdot,\cdot}\). (You may try
it yourself and find the process blocked by some terms that need additional information to push forward the computation.) Consider
\(G\pr{q,N} = \dfrac{F\pr{q,N}}{N}\). Then if \(G\pr{\cdot,\cdot}\) is
twice continuously differentiable (which indeed can be relaxed), we have
\[\ba
G\pr{q,N} & = \frac{F\pr{q,N}}{N} = p^D\pr{q}\cdot q - \frac{c\pr{Nq}}{N}\\
G_N\pr{q,N} & = -\frac{q\cdot c'\pr{Nq}\cdot N - c\pr{Nq}}{N^2}\\
G_{Nq}\pr{q,N} & = -qc''\pr{Nq}
\ea\]

If \(c\pr{\cdot}\) is concave, then \(G_{Nq}\pr{q,N}\ge 0\) and
\(G\pr{q,N}\) has increasing differences in \(\pr{q,N}\), so
\(q^*\pr{\cdot}\) weakly increases with \(N\). If \(c\pr{\cdot}\) is
convex, then \(G_{Nq}\pr{q,N}\le 0\) and \(G\pr{q,N}\) has increasing
differences in \(\pr{q,-N}\), so \(q^*\pr{\cdot}\) weakly decreases with
\(N\).
}


\section{Multivariate Comparative Statics}\label{multivariate-comparative-statics}

Consider a two-variable maximization problem:
\[\pr{x_1^*\pr{\theta},x_2^*\pr{\theta}} = \arg\max_{\pr{x_1,x_2}\in X\subset \RR^2}F\pr{x_1,x_2,\theta}.\]

If \(F\) merely has increasing differences in \(\pr{x_1,\theta}\),
that alone does not let us conclude that \(x_1^*\pr{\theta}\) is
weakly increasing — \emph{unless} \(x_2^*\) is independent of
\(\theta\). The reason: when \(\theta\) moves, it has a direct effect
on \(x_1\), but also an \emph{indirect} effect channeled through
\(x_2\) (since \(x_2^*\) may shift, which in turn changes the optimal
\(x_1\)). To sign the total effect we need a stronger structure on
\(F\) that controls both channels at once.

In the univariate case the strong set order used \(\min\brace{a,b}\)
and \(\max\brace{a,b}\). Extending this to vectors raises two
questions:

\begin{itemize}
\item
  How should we define ``min'' and ``max'' of two vectors?
\item
  Under that definition, will the result still lie in \(X\)?
\end{itemize}

For the first question, we introduce \emph{meet} and \emph{join} —
the componentwise minimum and maximum, respectively, which give the
greatest lower bound and the smallest upper bound of two vectors.

\defn{Meet; Join}{
For \(\x,\y\in\RR^n\), 
\begin{itemize}
    \item \textit{meet} of $\x$ and $\y$ is defined as
    \[ \x\meet \y = \pr{\min\brace{x_1,y_1},\min\brace{x_2,y_2},\cdots,\min\brace{x_n,y_n}}.\]
    \item \textit{join} of $\x$ and $\y$ is defined as
    \[\x \join \y = \pr{\max\brace{x_1,y_1},\max\brace{x_2,y_2},\cdots,\max\brace{x_n,y_n}}.\]
\end{itemize}

\hspace{2em} where \(\x\meet \y\) is called the greatest lower bound of \(\x\) and
\(\y\), and \(\x\join\y\) is called the smallest upper bound of \(\x\)
and \(\y\).
}

For the second question, we restrict attention to choice sets that are
\emph{closed under meet and join}. This structure is called a
\emph{sublattice}.

\defn{Sublattice}{
A set \(X\subset \RR^n\) is a
\emph{sublattice} if for any \(\x,\y\in X\), both \(\x\meet \y \in X\)
and \(\x\join \y\in X\).
}

\rmk{
\(\RR^n\) itself is a \emph{lattice}, so any set in
\(\RR^n\) is called a \emph{sublattice}.
}

\ex{
\begin{itemize}
\item
  \(X = X_1\times X_2\times \cdots\times X_n\), where \(X_i\subset \RR\),
  for \(i = 1,2,\cdots,n\).
\item
  \(X = \brace{\x\in \RR_+^n:\p\cdot\x\le m}\), where \(\p\gg\0\) is the
  price vector and \(m>0\) is the income. (NOT a sublattice)
\item
  \(X = \brace{\pr{x_1,x_2}\in \RR^2:x_2\le g\pr{x_1}, \text{ for }g\pr{\cdot} \text{ strictly increasing}}\).
\end{itemize}
}

The multivariate analogue of increasing differences is
\emph{supermodularity} — a single condition that captures
complementarity across all pairs of choice variables simultaneously.

\defn{Supermodularity}{
Let \(F:X\to \RR\), where
\(X\subset \RR^n\) is a sublattice. $F$ is \emph{supermodular} if
for any \(\x,\y\in X\),
\[F\pr{\x\meet\y}+F\pr{\x\join\y}\ge F\pr{\x}+F\pr{\y}.\]
}

\rmkb{
\begin{itemize}
\item
  Note when \(X = X_1\times X_2\subset \RR^2\), supermodularity is
  equivalent to increasing differences in \(\pr{x_1,x_2}\).
\item
  More generally, when
  \(X = X_1\times X_2\times \cdots\times X_n\subset \RR^n\),
  \textbf{supermodularity is equivalent to increasing differences in
  \(\pr{x_i,x_j}\) for all pairs of \(i\ne j\).}
\end{itemize}
}

Putting all together, we have the following multivariate
Topkis's Theorem.

\propp{Multivariate Topkis's Theorem}{
Let
\(F:X\times \Theta\to \RR\), where \(X\subset \RR^n\) is a sublattice and
\(\Theta\subset \RR\). Consider
\(\x^*\pr{\theta} = \arg\max_{\x\in X}F\pr{\x,\theta}\). If \(F\) is
supermodular, then for any \(\theta'>\theta\) and
\(\x\in \x^*\pr{\theta}\) and \(\x'\in \x^*\pr{\theta'}\), we have
\[\ba
\x\meet \x'&\in \x^*\pr{\theta},\\
\x\join \x'&\in \x^*\pr{\theta'}.
\ea\]
}{
Since \(X\) is a sublattice, \(\x\meet \x'\in X\)
and \(\x\join \x'\in X\). By revealed preference,
\[\bc
F\pr{\x,\theta}  \ge F\pr{\x\meet\x',\theta}\\
F\pr{\x',\theta'}  \ge F\pr{\x\join \x',\theta'}
\ec\too 
F\pr{\x,\theta}+F\pr{\x,\theta'}\ge F\pr{\x\meet \x',\theta}+ F\pr{\x\join\x',\theta'}\]
Since \(F\) is supermodular,
\[\ba
F\pr{\x,\theta}+F\pr{\x',\theta}& \le F\pr{\pr{\x,\theta}\meet \pr{\x',\theta'}}+F\pr{\pr{\x,\theta}\join \pr{\x',\theta'}}\\
& = F\pr{\x\meet \x',\theta}+F\pr{\x\join\x',\theta'}
\ea\]
Jointly, it must be that
\[F\pr{\x,\theta}+F\pr{\x,\theta'} = F\pr{\x\meet \x',\theta}+ F\pr{\x\join\x',\theta'}\]
Consequently, we have
\[\ba
F\pr{\x\meet\x',\theta} &= F\pr{\x,\theta}\\
F\pr{\x\join\x',\theta'} &= F\pr{\x,\theta'}
\ea\]

which indicates that \(\x\meet\x'\in \x^*\pr{\theta}\) and
\(\x\join\x'\in \x^*\pr{\theta'}\).
}

\ex{
Suppose a firm uses two inputs (\(L\) and \(K\)) to
produce a single output \(y\), and the production function is given by
\[y = f\pr{L,K}\]
\hspace{2em} where \(f\pr{\cdot,\cdot}\) is twice continuously differentiable.

\begin{enumerate}
\item
  Discuss how the optimal demand for capital \(K^*\pr{w,r,p}\) would be
  affected by an increase in the rental rate of the capital \(r\).
\item
  Discuss how the optimal demand for labor \(L^*\pr{w,r,p}\) would be
  affected by an increase in the rental rate of the capital \(r\).
\end{enumerate}
}

\sol{
\[\pi\pr{L,K,r} = pf\pr{L,K} - wL - rK.\]
Immediately, \(\pi_{Lr} = 0\) and \(\pi_{Kr} = -1\), \(\pi\) is supermodular in \(\pr{-K,r}\), so \(K^*\pr{w,r,p}\)
weakly decreases with \(r\). Next consider the indirect path of \(r\) to \(L\) 
through \(K\).

\begin{itemize}
\item
  If \(f_{LK}>0\), \(\pi\) is supermodular in \(\pr{-L,-K,r}\), so
  \(L^*\pr{w,r,p}\) weakly decreases with \(r\).
\item
  If \(f_{LK}<0\), \(\pi\) is supermodular in \(\pr{L,-K,r}\), so
  \(L^*\pr{w,r,p}\) weakly increases with \(r\).
\end{itemize}
}

\ex{
  Consider the other form of market power, a market with a single buyer and competitive supply. Suppose the inverse supply curse $p^S(q)\ge0$ is strictly increasing. Let $v(q)$ be the buyer's value for quantity $q$. Consequently, if there are $q$ units of transaction in the market, 
  \begin{itemize}
    \item The producer surplus is
    \[S(q) = \int_0^q\br{p^S(q)-p^S(t)}\d t.\]
    \item The buyer's gain is 
    \[B(q) = v(q)-p^S(q)q.\]
  \end{itemize}
  \begin{enumerate}
     \item Formulate both the buyer's and the social planner's optimization problems (no need to solve either of them).
    \item How does the single buyer's optimal quantity $q^B$ compare with that of the social planner $q^*$? Show your argument and explain intuitively.
  \end{enumerate}
}
\sol{
  \begin{enumerate}
    \item The buyer's optimization problem is
    \[\max_{q\ge 0}v(q)-p^S(q)q.\]
    The social planner's optimization problem is
    \[\max_{q\ge 0}v(q)-\int_0^qp^S(t)\d t.\]
    \item Consider the following optimization problem:
    \[
    \max_{q\ge 0} w(q,\lambda) = B(q)+\lambda S(q)
    \]
    This becomes the buyer's optimization problem when $\lambda = 0$, and the social planner's when $\lambda = 1$. If we could show how $q$ monotonically change with $\lambda$, then we are done with this question. This is the univariate comparative statics problem, so we check the increasing differences of $w(q,\lambda)$:
    \begin{equation*}
      \frac{\partial w(q,\lambda)}{\partial \lambda} = S(q)
    \end{equation*}
    Take any $q'\ge q\ge 0$. We have
    \begin{equation*}
      \begin{aligned}
        S(q')-S(q) & = \int_0^{q'}\br{p^S(q')-p^S(t)}\d t - \int_0^{q}\br{p^S(q)-p^S(t)}\d t \\
        & = \int_{q}^{q'}\br{p^S(q')-p^S(q)}\d t +\int_0^q\br{p^S(q')-p^S(q)}\d t \\
        & = \int_{q}^{q'}\br{p^S(q')-p^S(q)}\d t + \br{p^S(q')-p^S(q)}q\\
        & \ge 0
      \end{aligned}
    \end{equation*}
    where the last inequality holds because $p^S(\cdot)$ is upward sloping.

    Thus, $w(q,\lambda)$ has increasing differences in $(q,\lambda)$. From Topkis's Theorem, $q^*(\lambda)<q^*(\lambda')$.
  \end{enumerate}
}

Similar to increasing differences, supermodularity is
a \emph{cardinal} property, which is again too strong. Indeed, the
weaker requirement, \emph{quasi-supermodularity} serves our purpose.

\defn{Quasi-Supermodularity}{
Let \(F:X\to \RR\), where
\(X\subset \RR^n\) is a sublattice. \(F\) is
\emph{quasi-supermodular} if for any \(\x,\y\in X\),
\[\ba
F\pr{\x}\ge F\pr{\x\meet \y}&\too F\pr{\x\join \y}\ge F\pr{\y},\\
F\pr{\x}> F\pr{\x\meet \y}&\too F\pr{\x\join \y}> F\pr{\y}.
\ea\]
}

Under such extension, the analysis can be further extended beyond the
case of \(X\subset \RR^n\) and \(X\) forming a lattice structure.

